% !TEX root = ../main.tex

\chapter{预备知识}

本章介绍后续算法设计、系统实现与部署验证所需的共同基础，重点包括导航任务形式、视觉目标导航、扩散模型与扩散策略、NoMaD 的关键机制，以及轮式平台与足式平台在部署上的差异。由于本文研究同时涉及模型、系统与平台三个层面，因此预备知识不仅要给出概念定义，还要说明这些概念在本工作中的具体作用。

\section{导航与视觉导航基础}

\subsection{导航任务的基本定义}

机器人导航的目标是在给定观察、任务目标和运动约束的条件下，使机器人从当前状态移动到目标区域。一个完整导航系统通常包含环境感知、状态表示、目标表达、局部决策和运动执行五个环节\cite{thrun2005probabilistic,anderson2018evaluation}。在输出形式上，高层导航策略通常可以输出底层控制量、局部路径点或局部轨迹序列。本文采用第二类思路，即让高层模型输出局部 waypoint 轨迹，再由中层控制器完成执行映射。

如果把高层策略写成函数形式，则有
\begin{equation}
\pi : (\mathcal{O}_t, g) \mapsto \mathbf{A}_t ,
\end{equation}
其中 $\mathcal{O}_t$ 表示时刻 $t$ 的图像观察或观察序列，$g$ 表示任务目标，$\mathbf{A}_t$ 表示未来短时间内的局部轨迹。该表示天然适合与不同平台解耦，也是本文能够将高层视觉导航从轮式框架迁移到 Lite3 四足平台的重要原因。

\subsection{视觉导航与目标导航}

视觉导航是指以图像作为主要环境输入，并基于视觉信息完成导航决策的过程\cite{zhu2017targetdriven,anderson2018evaluation}。与严格依赖几何地图的传统方法相比，视觉导航更强调直接学习“观察到行为”的映射关系，因此更适合目标表达灵活、环境先验不足或快速部署的场景。

目标导航则进一步关注“目标如何给出”。常见形式包括 point-goal、image-goal、object-goal 和 language-goal 导航。本文主要研究 image-goal 场景：目标由一张图像或一组视觉节点图像给出。相较于绝对坐标，图像目标更符合真实应用，因为用户往往更容易提供“去这个画面对应的位置”，而非统一坐标系下的精确终点。

\subsection{拓扑导航与 topomap}

为了降低对精确几何地图的依赖，拓扑导航通常只保留若干代表性视觉节点及其连通关系\cite{savinov2018sptm,chaplot2020nts,shah2021ving}。在本文系统中，topomap 并不是在部署时由编码器在线生成的图结构，而是由预先采集到的一组候选节点图像构成。系统在运行时使用当前观察与这些候选节点的相对接近度进行筛选，再把选中的节点作为局部子目标输入 NoMaD。

设候选节点图像集合为
\begin{equation}
\mathcal{G}=\{I_g^{(1)}, I_g^{(2)}, \ldots, I_g^{(M)}\},
\end{equation}
距离预测头输出每个候选节点的相对距离估计 $\hat d_j$，则可通过
\begin{equation}
j^\ast = \arg\min_j \hat d_j
\end{equation}
选择当前更合适的局部子目标。由此可见，topomap 的作用是把长程任务拆解为一系列可被高层局部策略处理的短程子任务。

\begin{figure}[htbp]
    \centering
    \fbox{\rule{0pt}{52mm}\rule{0.82\linewidth}{0pt}}
    \caption{图位占位：导航任务输入输出与 topomap 关系示意图（待补）}
    \label{fig:prelim-nav-pipeline-placeholder}
\end{figure}

\section{扩散模型与扩散策略}

\subsection{扩散模型基本原理}

扩散模型通过“逐步加噪、逐步去噪”的方式学习复杂分布\cite{ddpm2020}。在前向过程里，数据样本 $x_0$ 被逐步扰动为高斯噪声：
\begin{equation}
q(x_t \mid x_{t-1}) = \mathcal{N}\!\left(\sqrt{1-\beta_t}x_{t-1}, \beta_t \mathbf{I}\right),
\end{equation}
其中 $\beta_t$ 为第 $t$ 步噪声调度系数。经过多步加噪后，样本会接近标准高斯分布。模型训练的目标是在反向过程中预测噪声，从而逐步恢复原始样本。常见训练目标写为
\begin{equation}
\mathcal{L}_{\text{diff}} = \mathbb{E}_{x_0,\epsilon,t}\left[\left\|\epsilon - \epsilon_\theta(x_t, t, c)\right\|_2^2\right],
\end{equation}
其中 $\epsilon$ 为真实噪声，$\epsilon_\theta$ 为模型预测噪声，$c$ 表示条件信息。

这种建模方式的优势在于，它不强迫模型在多峰分布中只输出一个均值解，而是允许模型通过去噪过程逐步生成多种合理候选。导航与控制任务中的未来轨迹本身就常常具有多模态特性，因此扩散建模天然适合此类问题。

\subsection{DDPM 与 DDIM}

DDPM 采用随机反向扩散过程生成样本，采样质量通常较好，但采样步数较多、推理开销较大\cite{ddpm2020}。DDIM 在保持训练目标不变的前提下，将反向过程改写为非马尔可夫且可确定性的更新形式，从而实现更少步数的快速采样\cite{ddim2020}。当取确定性采样形式时，其一步更新可写为
\begin{equation}
x_{t-1} = \sqrt{\bar\alpha_{t-1}}\hat x_0 + \sqrt{1-\bar\alpha_{t-1}}\,\epsilon_\theta(x_t,t,c),
\end{equation}
其中 $\hat x_0$ 由当前样本 $x_t$ 与预测噪声反推得到，$\bar\alpha_t$ 为累计噪声调度系数。

本文后续所说的“DDIM 推理加速”，就是在不重新训练 NoMaD 的前提下，用更少的去噪步数近似原始 DDPM 采样。它的核心问题不是“能不能更快”，而是“更快之后轨迹分布是否仍可接受”。

\subsection{Classifier-Free Guidance}

Classifier-Free Guidance（CFG）最初用于条件生成任务中的引导增强\cite{cfg2022}。其思想是在训练阶段同时学习有条件与无条件预测，在推理阶段通过两者线性组合增强条件约束。若条件分支噪声预测为 $\epsilon_{\text{cond}}$，无条件分支为 $\epsilon_{\text{uncond}}$，则引导后的噪声预测写为
\begin{equation}
\epsilon_{\text{cfg}} = (1+w)\epsilon_{\text{cond}} - w\epsilon_{\text{uncond}},
\end{equation}
其中 $w$ 为 guidance scale。当 $w=0$ 时退化为普通条件推理；当 $w>0$ 时，模型会更强地朝向条件约束；当 $w$ 过大时，也可能带来分布偏移与不稳定。

在本文中，CFG 的作用是增强轨迹对目标条件的牵引能力，即让生成的局部 waypoint 更明确地朝向目标图像或目标节点图像对应的方向。

\subsection{扩散策略与 Test-Time Scaling}

扩散策略将扩散模型从图像生成推广到动作序列或轨迹序列生成\cite{diffuser2022,diffusionpolicy2023}。在具身智能与机器人控制领域，Diffuser 和 Diffusion Policy 证明了扩散建模能够有效处理多步决策、长时序动作和多模态行为分布，因此被视为扩散模型在具身控制中的代表性应用。

在此基础上，近年来出现了测试时扩展（Test-Time Scaling, TTS）的思路\cite{ttsdiffusion2025,ttssearch2025}。其核心思想不是修改训练过程，而是在推理阶段使用额外计算预算生成多条候选轨迹，再通过打分函数进行选择。若生成 $N$ 条候选轨迹 $\{\mathbf{A}^{(i)}\}_{i=1}^N$，评分函数为 $s(\cdot)$，则最优轨迹可表示为
\begin{equation}
\mathbf{A}^{\ast} = \arg\max_{i\in\{1,\ldots,N\}} s\!\left(\mathbf{A}^{(i)} \mid \mathcal{O}_t, I_g\right).
\end{equation}

需要强调的是，TTS 与 DDIM、CFG 作用位置不同。DDIM 决定“怎么采样”，CFG 决定“采样时如何增强条件”，而 TTS 决定“采样出多条候选后如何选择”。因此三者虽然都属于推理时优化，但不是简单相加关系，而是具有先后层次的组合关系。

\subsection{扩散模型、ViT 与导航建模的关系}

Transformer 在视觉建模中的广泛应用，使得扩散模型与视觉编码方法的结合成为趋势。ViT 将图像切分为 patch token 并以自注意力机制进行建模\cite{vit2020}；DiT 则将扩散过程进一步与 Transformer 架构结合，说明扩散模型可以在强表征模型之上获得更高质量生成能力\cite{dit2022}。在导航领域，ViNT 使用 Transformer 处理历史观察与目标图像\cite{vint2023}，NoMaD 则在此基础上引入动作扩散策略\cite{nomad2023}。因此，本文讨论的视觉编码器升级、扩散采样优化和系统部署，其实都位于这一技术演化脉络之中。

\begin{figure}[htbp]
    \centering
    \fbox{\rule{0pt}{55mm}\rule{0.82\linewidth}{0pt}}
    \caption{图位占位：扩散模型、扩散策略与 NoMaD 关系示意图（待补）}
    \label{fig:prelim-diffusion-placeholder}
\end{figure}

\section{NoMaD 关键机制}

\subsection{NoMaD 的基本结构}

NoMaD 由视觉编码器、距离预测头和扩散策略头三部分构成\cite{nomad2023}。视觉编码器负责把历史观察与目标图像编码为统一条件特征；距离预测头估计当前观察与目标之间的相对距离；扩散策略头则在该条件特征上生成未来局部 waypoint 轨迹。若记局部轨迹为
\begin{equation}
\mathbf{A} = [\mathbf{a}_1,\mathbf{a}_2,\ldots,\mathbf{a}_H], \qquad \mathbf{a}_i=(\Delta x_i,\Delta y_i),
\end{equation}
则 NoMaD 学习的是条件分布
\begin{equation}
p(\mathbf{A}\mid \mathcal{O}_t, I_g).
\end{equation}

\subsection{距离预测头与视觉编码器的直接耦合}

在当前项目实现中，距离预测头与扩散策略头共享同一条视觉条件空间。也就是说，视觉编码器一旦更换，距离预测头看到的输入分布和扩散策略头看到的输入分布会同时改变。因此，视觉编码器并非独立前端，距离预测头也不是与编码器无关的附属模块。本文后续视觉编码器实验之所以必须重新训练距离分支，原因就在于这种直接耦合结构。

\subsection{goal mask 与导航/探索统一建模}

NoMaD 的另一个关键机制是 goal mask。训练时，模型以一定概率遮蔽目标 token，使其同时经历“有目标条件”和“无显式目标条件”两种情况\cite{nomad2023}。这并不意味着模型在无目标条件时进行随机游走，而是意味着它会学习在当前视觉上下文中生成连续、可执行且符合数据轨迹分布的探索性行为。正因为如此，NoMaD 才能在部署时同时支持导航模式与探索模式。

\section{轮式平台、足式平台与 Lite3}

\subsection{轮式与足式平台的部署差异}

轮式平台姿态变化小、视觉输入稳定，与速度控制器的接口也较简单；足式平台则具有更强的地形适应性，但会带来更明显的机体起伏、更严格的实时性约束和更复杂的控制接口\cite{hwangbo2019anymal,lee2020learning}。这意味着，同一高层轨迹若能在轮式平台上顺利执行，并不代表它能直接在四足平台上稳定工作。

\begin{table}[htbp]
    \centering
    \caption{轮式平台与足式平台在视觉导航中的主要差异}
    \label{tab:wheel-leg-compare}
    \begin{tabular}{lll}
        \toprule
        维度 & 轮式平台 & 足式平台 \\
        \midrule
        机体姿态变化 & 较小 & 明显起伏与抖动 \\
        视觉输入稳定性 & 较好 & 易受步态扰动影响 \\
        高层到低层接口 & 较直接 & 更依赖中层解耦 \\
        实时性约束 & 较高 & 更高 \\
        安全机制需求 & 相对简单 & 更严格 \\
        \bottomrule
    \end{tabular}
\end{table}

\subsection{云深处 Lite3 平台与关键参数}

本文采用的主要足式平台为云深处 Lite3。结合项目配置与部署代码，其在本文中的系统定位是“由高层视觉导航策略驱动的四足运动系统”，而不是端到端关节控制平台。对应控制链路为
\begin{equation}
\text{视觉观察} \rightarrow \text{高层 NoMaD 推理} \rightarrow \text{中层控制映射} \rightarrow \text{底层运动执行}.
\end{equation}

与本文实现直接相关的关键参数如表~\ref{tab:lite3-params} 所示。
\begin{table}[htbp]
    \centering
    \caption{本文中与 Lite3 相关的关键平台参数}
    \label{tab:lite3-params}
    \begin{tabular}{lll}
        \toprule
        参数类别 & 参数项 & 数值或说明 \\
        \midrule
        平台类型 & 机器人形态 & 四足平台 \\
        运动结构 & 关节数量 & 12 关节执行系统 \\
        速度接口 & 前后速度范围 & $[-1.0,1.0]\,\text{m/s}$ \\
        速度接口 & 左右速度范围 & $[-0.5,0.5]\,\text{m/s}$ \\
        速度接口 & 偏航角速度范围 & $[-1.5,1.5]\,\text{rad/s}$ \\
        部署限幅 & 本文最大线速度 & $0.2\,\text{m/s}$ \\
        部署限幅 & 本文最大角速度 & $0.4\,\text{rad/s}$ \\
        感知频率 & observation frame rate & $4\,\text{Hz}$ \\
        通信安全 & 速度超时保护 & $250\,\text{ms}$ 自动停止 \\
        \bottomrule
    \end{tabular}
\end{table}

这些参数并非纯说明信息，而是会反向约束系统设计。较低的高层观察频率要求单次推理延迟足够低；保守的速度限幅要求轨迹输出足够平滑；超时保护要求系统必须具备显式的安全停机和恢复机制。因此，本文后续才会在系统设计中引入统一推理模块、状态机和导航主机。

\begin{figure}[htbp]
    \centering
    \fbox{\rule{0pt}{55mm}\rule{0.82\linewidth}{0pt}}
    \caption{图位占位：Lite3 平台与控制链路示意图（待补）}
    \label{fig:prelim-lite3-placeholder}
\end{figure}

\section{本章小结}

本章从任务形式、模型原理和平台约束三个层面为后续章节建立了共同语言。首先说明了导航、视觉导航、目标导航与 topomap 的基本作用；随后详细介绍了扩散模型、DDIM、CFG、扩散策略和 TTS 的原理、公式与引用脉络；接着总结了 NoMaD 的结构、goal mask 机制以及视觉编码器与距离预测头的直接耦合关系；最后说明了轮式与足式平台的差异，以及 Lite3 在本文中的平台定位与关键参数。基于这些预备知识，下一章将进一步进入 NoMaD 的算法设计与离线实验分析。
